% !TEX root = ../main.tex
\section{Category I: Prompting-Based Methods}\label{sec:prompting}

\begin{summarybox}[s1]{Paradigm Summary \hfill \promptbadge}
Prompting-based methods keep all LLM weights \textbf{frozen} and perform failure attribution purely through inference-time scaffolding: log partitioning, failure taxonomies, programmatic/semantic constraint checking, dependency and causal graph construction, or iterative multi-agent adjudication. They require \emph{no model training} and immediately inherit gains from stronger foundation models, but face fundamental tensions between context length, inference token cost, and vulnerability to late-stage symptom bias.
\end{summarybox}

Across the literature, prompting-based attribution has rapidly progressed through four increasingly structured architectural waves.

\subsection{Direct and Sequential Log-Scanning Judges}
The earliest wave feeds raw or minimally segmented execution logs directly to an LLM judge.

\paragraph{Who\&When Baselines~\citep{zhang2025who}}
Formulating automated MAS failure attribution alongside the $184$-trace Who\&When benchmark ($126$ CaptainAgent Algorithm-Generated and $58$ Magentic-One Hand-Crafted logs), Zhang et al.\ establish three canonical prompting protocols:
(i)~\textbf{All-at-Once}, which feeds the entire trajectory in a single prompt to predict $(i^*, t^*)$---giving global visibility at low call count ($O(1)$) but suffering severe long-context dilution;
(ii)~\textbf{Step-by-Step}, which incrementally reveals prefixes $\tau_{\le t}$ and halts at the first detected mistake---sharpening local focus at $O(T)$ cost but lacking downstream outcome awareness; and
(iii)~\textbf{Binary Search}, which recursively halves the trace in $O(\log T)$ calls---improving scalability but failing when deciding which half contains a subtle upstream root cause. Even frontier reasoning models (o1, DeepSeek-R1) achieved at most $53.5\%$ agent accuracy and $14.2\%$ step accuracy, collapsing near $0\%$ on the longest logs.

\paragraph{Taxonomy-Guided Judges: MAST~\citep{cemri2025mast}, TRAIL~\citep{deshpande2025trail}, and AgentFail~\citep{ma2025agentfail}}
Several foundational works pair raw log judging with empirically grounded failure taxonomies. \textbf{MAST} derives $14$ failure modes across $3$ categories (System Design, Inter-Agent Misalignment, Task Verification; human $\kappa=0.88$) from $1{,}642$ traces and builds a calibrated OpenAI o1 LLM-as-a-Judge pipeline ($\kappa=0.77$--$0.79$ with humans). \textbf{TRAIL} targets OpenTelemetry-formatted single- and multi-agent traces ($148$ traces, $841$ span-level errors across Reasoning, Execution, and Planning), showing that even Gemini-2.5-Pro scores only $11\%$ joint localization--classification accuracy. \textbf{AgentFail} studies $307$ failure logs from low-code workflow platforms (Dify, Coze, n8n, AutoGen), demonstrating that injecting its platform/workflow/agent taxonomy into All-at-Once, Step-by-Step, and Binary Search boosts root-cause accuracy by $15$--$20$ percentage points (from $8.3\%$--$13.0\%$ to $24.1\%$--$33.6\%$), yet leaving two-thirds of cases unsolved.

\paragraph{Full-Observability Auditing and Two-Stage Scoping: TraceElephant~\citep{chen2026traceelephant} and SDBL~\citep{sun2026sdbl}\,\webtag}
\textbf{TraceElephant} ($220$ fully observable, reproducible traces across CaptainAgent, Magentic-One, and mini-SWE-agent) reveals that supplying full step inputs, agent prompts, and architecture specs instead of output-only logs lifts LLM attribution accuracy by up to $76\%$, and that an agentic tool-equipped inspector outperforms static scanning. \textbf{SDBL} (\emph{Scope Delineation Before Localization}, AAAI 2026) directly addresses context overload by decoupling attribution into two prompting stages: first narrowing a long trajectory to a compact suspicious candidate window (via stepwise or expertise-assisted scoping), then performing fine-grained localization only inside the scoped subset ($68.5\%$ HC / $63.5\%$ AG agent; $27.8\%$ HC / $39.7\%$ AG step on Who\&When w/o G).

\subsection{Structured Context, Constraint, Schema, and Claim-Auditing Pipelines}
To overcome brittle raw-text judging, a second wave transforms traces into structured intermediate artifacts---compressed hierarchies, modular checks, explicit constraints, cached error schemata, or claim ledgers---before final adjudication.

\paragraph{ECHO~\citep{banerjee2025echo}}
ECHO (\emph{Error attribution through Contextual Hierarchy and Objective consensus}) tackles long-context degradation via three decoupled mechanisms: (1) a \textbf{four-layer regex-based hierarchical context} around each target step---$L_1$ immediate neighbors ($\tau_{i\pm 1}$, verbatim), $L_2$ local ($\tau_{i\pm 2,3}$, conclusions/transitions), $L_3$ distant ($\tau_{i\pm 4..6}$, state/error summaries), and $L_4$ global milestones; (2) a \textbf{panel of 6 specialized analyst personas} (Conservative, Liberal, Detail-Focused, Pattern-Focused, Skeptical, General; $3$ sampled per trace at temperatures $0.3$--$0.9$) outputting step likelihoods and alternative hypotheses; and (3) \textbf{confidence-weighted consensus voting} ($\delta=0.3$). With Claude-3.5-Sonnet, ECHO reaches $68.4\%$ HC / $68.8\%$ AG agent accuracy and $28.1\%$ HC / $28.8\%$ AG step accuracy on Who\&When; switching regex compression to LLM extraction raises HC step accuracy to $43.8\%$ but triples token cost ($166\text{K}$ vs.\ $54\text{K}$ tokens).

\paragraph{AgentDebug~\citep{zhu2026agentdebug}}
Targeting single-agent rollouts decomposed into four cognitive modules (\emph{Memory, Reflection, Planning, Action}) on AgentErrorBench (ALFWorld, GAIA, WebShop), AgentDebug executes a three-stage prompting loop: (1) fine-grained per-step, per-module mapping onto the modular \emph{AgentErrorTaxonomy}; (2) LLM critical-error detection selecting the earliest root-cause step $t^*$; and (3) iterative re-rollout from $t^*$ with targeted corrective feedback. On AgentErrorBench with GPT-4.1, AgentDebug achieves $42\%$ step accuracy and $32\%$ strict all-correct accuracy ($+17\%$ step and $+24\%$ all-correct over baselines), driving up to $+26\%$ relative downstream task recovery.

\paragraph{AgentRx~\citep{barke2026agentrx}}
Inspired by software design-by-contract, AgentRx diagnoses the \emph{first unrecoverable failure} across API workflows ($\tau$-bench), incident management (Flash), and web/file MAS (Magentic-One). Given domain policies and tool schemas, AgentRx automatically synthesizes both \textbf{programmatic constraints} and \textbf{semantic natural-language guarded constraints}, evaluates each constraint step-by-step to produce an auditable \textbf{validation log} $\mathcal{V} = \{(k, C, e)\}$ of violations with grounded evidence, and feeds $\mathcal{V}$ plus a grounded-theory $9$-category checklist to an LLM judge. Across the three domains (using GPT-5), AgentRx delivers $+23.6\%$ absolute step-localization and $+22.9\%$ failure-category improvements over direct baselines.

\paragraph{AgentScope~\citep{bi2026agentscope}}
AgentScope advances a \textbf{neuro-symbolic} diagnostic pipeline: (1) instrumentation and selective LLM parsing convert raw traces into a \textbf{Reasoning-Action Graph (ReAG)} with aligned semantic fields (role, step purpose, actions, missing-info signals); (2) a library of \textbf{neural invariants}---formal failure-mode predicates backed by targeted JSON-returning LLM functions spanning $10$ failure types across Reasoning, Control-flow, and Action---scans the ReAG to retain multiple evidence-backed candidate hypotheses; and (3) a decisive-error judge selects the single fault whose downstream impact best explains task failure. AgentScope achieves $25.4\%$--$77.8\%$ step accuracy on Who\&When-AG, $22.4\%$--$34.5\%$ on HC, and $28.4\%$--$54.1\%$ on the newly introduced $303$-trace AgentErrata benchmark.

\paragraph{CORRECT~\citep{yu2026correct}}
Observing that $>80\%$ of failed Who\&When traces exhibit pairwise BERT cosine similarity $>0.8$ yet naive few-shot ICL degrades accuracy due to $32\text{K}+$ context bloat, CORRECT introduces a training-free \textbf{online error-schema cache}. From seed annotated failures, an LLM distills compact structured schemata $S = \langle \Sigma, C, H \rangle$ containing \emph{Error Signatures} $\Sigma$, \emph{Error Context Analysis} $C$, and actionable \emph{Detection Heuristics} $H$. At inference, semantically relevant schemata are retrieved in sub-seconds to guide targeted localization, yielding up to $+19.8\%$ accuracy gains on Who\&When and $+28.7\%$ average gains across seven MAS applications on CORRECT-Error ($2{,}000+$ traces) with zero LLM retraining.

\paragraph{Domain-Specialized Auditing: DRIFT~\citep{wang2026drift}, SearchAuditor~\citep{liang2026searchauditor}, and TrajDebug~\citep{qi2026trajdebug}\,\webtag}
Long-horizon search and deep-research agents rarely crash with exceptions; instead, they commit early to unsupported facts that look fluent downstream.
\textbf{DRIFT} reframes span-level localization (on its $1{,}000$-instance TELBench) around \textbf{claim-centric auditing}: (A)~\emph{Claim Keeper} builds a global claim ledger $\mathcal{L} = \{(a_k, i_k, b_k, U_k, \tau_k, \sigma_k)\}$ separating benign exploratory queries from settled commitments; (B)~\emph{Support Seeker} verifies trajectory evidence (\code{DIRECT/WEAK/...}); and (C)~\emph{Dependency Tracer} propagates unsupported commitments forward/backward to flag consequential error spans, boosting F1 and first-error accuracy by up to $+30$~pp.
\textbf{SearchAuditor} (on SearchAuditBench: $1{,}243$ trajectories averaging $73.1$ messages / $65.1\text{K}$ tokens) runs \textbf{three complementary parallel audit branches}, adjudicates their findings over a compact evidence-grounded trajectory view, and synthesizes rubric-graded repair directives---lifting GPT-5.5 end-to-end pass rate from $26.6\%$ to $32.3\%$ and fixing $17.4\%$ of Kimi-K2.6's live failures on LiveBrowseComp.
\textbf{TrajDebug}~\citep{qi2026trajdebug} combines multi-granularity history compression with \textbf{error-lifecycle state tracking}, explicitly distinguishing local errors that were subsequently resolved from persistent faults that drove final failure.

\subsection{Causal-Reasoning, Dependency-Graph, and Propagation-Guided Prompting}
A third---and currently highest-performing---prompting wave moves beyond flat logs and independent step checks by prompting LLMs to reconstruct explicit \textbf{causal, information-dependency, or error-propagation graphs} and verify counterfactual decisiveness.

\paragraph{A2P Scaffolding~\citep{weng2025a2p}}
A2P (\emph{Abduct--Act--Predict}, co-authored by DeepScientist) bridges the counterfactual inference gap within a \emph{single} LLM pass by imposing Pearl's three-step causal ladder alongside explicit contextual step numbering: (1)~\textbf{Abduction} infers latent root causes (assumptions, knowledge gaps) behind observed actions; (2)~\textbf{Action} specifies a minimal corrective intervention; and (3)~\textbf{Prediction} mentally simulates the post-intervention trajectory to verify task resolution. On Who\&When, A2P boosts step accuracy $2.85\times$ on AG ($47.46\%$) and $2.43\times$ on HC ($29.31\%$) with only $\sim 25\%$ latency overhead; ablations reveal that explicit step numbering alone accounts for $+29.7$~pp.

\paragraph{CHIEF~\citep{wang2026chief}}
CHIEF transforms flat logs into a two-level \textbf{Hierarchical Causal Graph (HCG)} $G=(V_{\text{sub}}\cup V_{\text{agt}}, E)$ in three phases: (1)~RAG-guided task decomposition (retrieving exemplar plans from GAIA/AssistantBench) refined by trajectory-aligned reflection yields contiguous subtask nodes and intra-subtask agent DAGs; (2)~\textbf{Hierarchical Oracle-Guided Backtracking} synthesizes virtual oracles to prune irrelevant subtasks top-down; and (3)~\textbf{Progressive Counterfactual Causal Screening} tests four counterfactual patterns (local, planning-control, data-flow, deviation-aware) to separate root causes from propagated symptoms. With DeepSeek-V3.2 (thinking), CHIEF attains $77.6\%/72.4\%$ HC and $76.8\%/68.8\%$ AG agent accuracy, and $29.3\%$ HC / $52.0\%/45.6\%$ AG step accuracy (w/ G $\mid$ w/o G) at $2.5$--$3\times$ All-at-Once token cost ($6$--$8\times$ cheaper than replay-based FAMAS).

\paragraph{EMFA~\citep{liao2026emfa}}
Noting that prior graph methods miss persistent unresolved interaction loops, EMFA (\emph{Error-Propagation Modeling for Failure Attribution}) executes five stages without requiring task ground truth: (1)~extracts a query-only success specification $C_{\text{success}}$ to anchor the terminal failure symptom $x_{\text{final}}$; (2)~constructs a two-level step/subtask dependency structure $\mathcal{H}$; (3)~explicitly models both \textbf{cascading propagation chains} and \textbf{persistent interaction loops}; (4)~\textbf{dual-branch candidate screening} extracts top candidates from both branches; and (5)~\textbf{counterfactual decisive-error verification} selects $(a^*, i^*)$. With GPT-4o (w/o G), EMFA sets strong step accuracy on Who\&When: $32.76\%$ HC ($+3.45$~pp over CHIEF) and $50.00\%$ AG ($+4.40$~pp), with $74.14\%$ HC agent accuracy.

\paragraph{DeFA~\citep{anonymous2027defa}}
DeFA (\emph{Dependency-Guided Failure Attribution}) decouples \textbf{protocol relations} (framework-induced conventions such as delegation or tool-call/return pairs) from \textbf{semantic dependencies} (implicit cross-step content reuse) to build an \textbf{Event Dependency Graph (EDG)}. To scale to long and multimodal (image/video) traces, DeFA partitions trajectories into segments, supplying each segment with full local detail plus chronological summaries of all other segments for candidate event recall. It then constructs a \textbf{Failure Propagation Graph (FPG)} tracing upstream sources and downstream amplification/repair, before an LLM adjudicator selects $(t^*, a^*, y^*)$. On Who\&When (combined $184$ traces, w/o G), DeFA reaches $72.3\%$ agent and $52.7\%$ exact step accuracy with GLM-5.2 ($80.3\%$ step on Who\&When Pro text), and feeding DeFA diagnoses into Trace2Skill yields $+6$--$15$~pp downstream task gains.

\paragraph{BLADE~\citep{guo2026blade}}
BLADE (\emph{Blame Localization via Audited Dependency Evidence}) makes a sharp conceptual distinction: \emph{tracing where information came from is not the same as auditing how it was used}. Using the structured probabilistic judgment interface Jev, BLADE (1)~computes probabilities over four directed relations (\code{depends\_on}, \code{repeats}, \code{contradicts}, \code{repairs}) and retains edges $\ge \theta=0.3$; (2)~plans selective \textbf{contract auditing} over only the top-$K=2$ highest-probability sources per event ($13\%$--$23\%$ of all pairs), checking whether each receiving decision faithfully respected its source and task requirement spans (\emph{error introduction vs.\ propagation vs.\ repair}); and (3)~performs \textbf{defeasible global diagnosis} over the full blame graph. Across $1{,}228$ trajectories spanning five Who\&When and Who\&When Pro subsets (w/o G), BLADE beats All-at-Once by $2.2$--$11.1$~pp (macro average $66.8\%$ agent, $57.6\%$ step; $54.0\%$ Who\&When-AG step).

\paragraph{Stateful Memory and Emerging Graph-Prompting Works}
\textbf{MemTrace}~\citep{deng2026memtrace} extends dependency-graph attribution to stateful non-parametric memory systems (RAG, Mem0, EverMemOS) by instrumenting source code to emit a bipartite variable--operation execution graph $G=(V,O,E)$, then diagnosing extraction, update, retrieval, deletion, and response errors via graph-walking or regex-based operation search (\code{MemTrace-OBS})---enabling closed-loop prompt optimization ($+7.62$~pp end-task accuracy).
Recent web-sourced methods reinforce this graph-guided trend:
\textbf{FALAT}~\citep{rafi2026falat}\,\webtag frames attribution as hierarchical dependency-guided search using task-conditioned priors, typed dependencies to prune downstream carriers, and counterfactual sufficiency checks ($73.0\%$ HC agent, $29.1\%$ HC / $38.0\%$ AG step);
\textbf{Adaptive Influence Graphs (AIG)}~\citep{bakish2026aig}\,\webtag dynamically defines trace-specific node boundaries and non-adjacent inheritance edges for agent-directed traversal;
\textbf{EDGE}~\citep{hou2026edge}\,\webtag targets \emph{multi-error} attribution on TRAIL and MAST by validating an Error Dependency Graph via counterfactual rollouts to steer a two-stage LLM detector;
\textbf{TRIAGE}~\citep{wang2026triage}\,\webtag ranks candidate steps by severity and structural impact; and
\textbf{CockpitHAT}~\citep{wang2026cockpithat}\,\webtag applies dependency-graph hierarchical attribution to embodied multi-agent cockpits.

\subsection{Iterative Multi-Agent Adjudication and Continual Search}
Rather than relying on a fixed single-pass graph or static prompt pipeline, a fourth prompting family treats debugging as an \textbf{iterative, tool-using investigation} by specialized diagnostic agents.

\paragraph{RAFFLES~\citep{zhu2026raffles}}
RAFFLES formalizes the distinction between \emph{trivial} and \emph{critical} step faults and introduces an iterative \textbf{Judge--Evaluator} architecture: a central Judge proposes candidate faults and rationales, while a panel of specialized Evaluators independently critiques both the fault selection and the Judge's reasoning with natural-language confidence scores; the Judge reflects on these critiques over $K$ refinement iterations. On Who\&When with Claude-Sonnet-4, RAFFLES achieves $51.59\%$ step accuracy on AG and $20.69\%$--$27.59\%$ on HC, while exceeding $83\%$ on mathematical reasoning (ReasonEval).

\paragraph{ErrorProbe~\citep{li2026errorprobe}}
ErrorProbe combines taxonomy cues, backward tracing, multi-agent tool-grounded verification, and self-improving memory without updating LLM weights: (1)~a lightweight MAST detector tags local anomalies as heuristic priors; (2)~symptom-driven backward tracing prunes long histories to causally relevant turns; (3)~a three-agent diagnostic team (\textbf{Strategist, Investigator, Arbiter}) validates hypotheses using executable tools (e.g., code sandboxes); and (4)~a \textbf{verified-before-write episodic memory} commits reusable error/patch patterns \emph{only} when confirmed by executable evidence, preventing memory drift and enabling zero-training cross-domain transfer (e.g., KodCode $\to$ TracerTraj, reaching $73.2\%$ agent / $39.4\%$ step on TracerTraj and $60.3\%$ agent on Who\&When-Algo with Claude-3.7-Sonnet).

\paragraph{Continual Search (MegaRCA)~\citep{raj2026megarca}\,\webtag}
Arguing that one-shot LLM judges lock prematurely onto salient late symptoms in massive logs, Raj et al.\ cast long-horizon root-cause attribution as an iterative \textbf{search problem}, introducing Continual Search to repeatedly nudge diagnostic models toward unexamined evidence regions across long trajectories.
